\chapter*{Abstract}
\addcontentsline{toc}{chapter}{Abstract}

\begingroup
\setstretch{1.0}
\rmfamily\fontsize{11}{11}\itshape\selectfont
\setlength{\parskip}{6pt}

This investigation evaluates the structural performance of sustainable polymer 
al\-tern\-at\-ives, specifically Recycled Poly\-car\-bon\-ate (rPC) and a 
Bamboo-HDPE 30\% wt composite, against standard Virgin High-Density 
Polyethylene (HDPE) and ABS safety helmets. Designed to meet the EN 397 industrial 
safety standard, construction helmets must restrict the transmitted force 
to a 5.0 kg headform to a maximum of 5.0 kN under a 49 Joules impact. 
Applying LS-DYNA explicit finite element analysis (FEA), the simulations 
modeled an upward relative impact velocity of 4.43 m/s based on the 
principle of Galilean Invariance.

A five-phase analytical methodology concluded with an advanced 30$^\circ$ 
oblique impact simulation to quantify rotational kinematics. 
Thermodynamic validation confirmed computational stability with 
exactly 0\% hourglass energy across all material models. Comparison of 
raw nodal acceleration data to filtered SAE J211 CFC 60 signals isolated the 
true macroscopic structural deceleration from numerical contact noise.

Results identified two primary energy-absorption mechanisms: highly ductile 
polymers (HDPE) absorbed impact energy through internal strain yielding, 
while high-modulus materials (ABS, rPC) used load distribution due to 
increased flexural stiffness to distribute impact forces across a wider surface area of the 
EPS liner. The Recycled PC shell demonstrated superior 
performance in oblique loading, achieving a 74\% reduction in peak 
rotational acceleration compared to Virgin HDPE. This is attributed to 
reduced tangential frictional interaction, where the higher flexural modulus 
limits localized surface deformation and reduces tangential torque compared 
to the plastic deformation observed in HDPE.

The findings confirm that sustainable material alternatives provide 
sufficient structural protection and meet safety-critical regulatory requirements while 
significantly reducing the environmental footprint of PPE production.\par
\endgroup

\vspace{12pt}
\noindent{\rmfamily\fontsize{12}{13.8}\itshape\selectfont \textbf{Keywords:} Explicit Dynamics, EN 397, FEA, Sustainable Polymers, Impact Analysis, SAE J211, Oblique Impact, Rotational Kinematics\par}
